\documentclass{article}
\usepackage{amsmath}
\usepackage{amsthm}

\newtheorem{theorem}{Problem}

\begin{document}

\begin{theorem}
Find the largest value of $x$ that satisfies the equation $|5x-1|=x+3$.
\end{theorem}

\begin{proof}
Step 0: The first thing I notice is that this equation involves the absolute value function, which means it can have two possible cases: either the expression inside the absolute value is positive or negative.
Step 1: If it is positive, then the absolute value does not change anything, and I can rewrite the equation as $5x-1=x+3$.
Step 2: Solving for $x$, I get $4x=4$ and $x=1$.
Step 3: However, I need to check if this value actually satisfies the original equation, because it might not work for the other case of the absolute value.
Step 4: Plugging in $x=1$, I get $|5(1)-1|=1+3$, which simplifies to $|4|=4$, which is true.
Step 5: So $x=1$ is a valid solution.
Step 6: Now I need to consider the other case, when the expression inside the absolute value is negative.
Step 7: This means that $5x-1<0$, or equivalently, $x<\frac{1}{5}$.
Step 8: In this case, the absolute value function changes the sign of the expression, so I can rewrite the equation as $-5x+1=x+3$.
Step 9: Solving for $x$, I get $-6x=2$ and $x=-\frac{1}{3}$.
Step 10: Again, I need to check if this value satisfies the original equation, by plugging it in and simplifying.
Step 11: Plugging in $x=-\frac{1}{3}$, I get $|5(-\frac{1}{3})-1|=-\frac{1}{3}+3$, which simplifies to $|\frac{8}{3}|=\frac{8}{3}$, which seems to be correct, but let's re-evaluate.
Step 12: Actually, the calculation is incorrect because $5(-\frac{1}{3})-1 = -\frac{5}{3} - 1 = -\frac{8}{3}$, so $|-\frac{8}{3}|=\frac{8}{3}$. Therefore, the simplification is true but we need to ensure alignment with inequality.
Step 13: However, re-evaluating shows $x=-\frac{1}{3}$ indeed satisfies the condition but lies outside $x < \frac{1}{5}$. Therefore, it's invalid under the inequality domain considered.
Step 14: So $x=-\frac{1}{3}$ is not in the range expected by our assumption on $x$.
Step 15: Therefore $x=1$ remains the only valid and largest value under true evaluation. It is within $x>\frac{1}{5}$ domain.
Step 16: # Answer\n1
\end{proof}

\end{document}
